\properlane{restored-vineyard}{The wounded vineyard made one}
The vineyard's wound calls forth more than an account of failure. It calls forth a plea that God visit what he planted. This second reading follows that plea toward Christ and toward the common life of his body. Augustine and Bellarmine understand the vineyard's restoration through Christ without making it a wholly unrelated planting. Chrysostom's Eucharistic exposition then asks whether people who receive the same Christ will live as members of one another. The governing question is how a wounded people is renewed, not how it can congratulate itself on surviving.

\subsection{The voice inside the broken enclosure}
Isaiah speaks the owner's accusation; the psalm speaks from within the damaged vineyard. The change of voice matters. In the reading, the hearer is brought to recognize the seriousness of unjust fruit. In the response, the assembly addresses the planter and asks him to act. The same people can need both judgment and mercy. Yet these voices do not authorize a universal diagnosis that every sufferer's wound measures personal guilt. The texts give particular accusations and particular laments, not a formula assigning blame to every injured person.

Mordecai's Entrance prayer adds corporate intercession before the vineyard is named. The threatened people appeal to a sovereignty greater than the court's power. Their hope does not depend on having become socially invulnerable. It depends on the Creator's ability to hear and preserve. The Collect likewise names need before merit and divine abundance beyond the requester's reach. Together these opening elements establish that restoration can be sought by those who do not command its means.

The psalm remembers Egypt, planting and expansion before describing broken defenses. Memory keeps the damaged community from understanding itself solely through the latest disaster. The ruin is real, but it does not abolish the identity of the one who planted or the history of his gift. Petition asks him to return to his work. Its verbs---look, visit, give life---express dependence upon an attention the people cannot manufacture for themselves.\footnote{Psalm 79:9,12--16,19--20, Douay--Rheims; Augustine, \work{Exposition of Psalm} 79, §§6--11; Bellarmine, \work{Commentary on the Book of Psalms}, Psalm 79.}

\subsection{Augustine: the same vine perfected in Christ}
Augustine moves between the vine's Israelite beginning and its fulfillment in Christ. He first acknowledges the people brought from Egypt and the territory they occupied. When he considers the vine's eventual reach, he reads its expansion through the way opened in Christ. The poem's images thus carry both the remembered planting and a perfection extending beyond its earlier borders. His exposition does not require the planter to discard one people and invent an unrelated origin.\footnote{Augustine, \work{Exposition of Psalm} 79, §6, NPNF1, vol. 8, heading Psalm 80.}

At the request that God perfect what his right hand planted, Augustine makes that continuity explicit. He corrects the suggestion of a second vineyard into the same vineyard, and connects the seed of Abraham with the promise of blessing to the nations. Gentile incorporation is described through grafting into an existing root. The Son of Man is the place of perfection: Christ is the foundation upon which the building stands. The vine's renewal therefore has a Christological center, rather than resting on an undifferentiated confidence in communal endurance.\footnote{Augustine, Psalm 79, §9.}

The Gospel's violence does not defeat this account. Augustine recalls the tenants and the heir in his psalm exposition. The heir cast out of the vineyard nevertheless possesses it more fully. Rejection does not strip the Son of his relation to the Father's work. That paradox corresponds to Matthew's movement from murdered heir to rejected cornerstone: the act intended to eliminate him becomes unable to control the outcome. The people's renewal is founded upon the one whom grasping stewards attempted to remove.\footnote{Augustine, Psalm 79, §7; Matthew 21:38--43.}

This is a Christian interpretation of the psalm's prayer, and Augustine gives it in the language of fulfillment. The literal vine remains the people brought from Egypt. Christological perfection deepens the reading of their prayer; it does not make their history an empty shell. Nor can the nations receiving the promise claim an independent right to the root into which they are grafted. The image itself requires gratitude for an inheritance that precedes their arrival.

\subsection{Bellarmine: reform, not total destruction}
Bellarmine gives an independent argument at the same psalmic turn. In his explanation of the plea for visitation, he distinguishes a reform of the vineyard from its complete destruction. He names the apostles as belonging to Israel and recalls the many who respond to Peter's preaching. The Gentiles enter through grafting, as in Paul's olive-tree image. These particulars matter because they make continuity concrete: the renewed people is not produced by erasing its Jewish members and beginning again elsewhere.\footnote{Bellarmine, \work{Commentary on the Book of Psalms}, Psalm 79, explanation at vv. 14--18 in the retained O'Sullivan translation (1866). The translation is abridged; the lemma numbering in the exposition does not supply a new numbering for the appointed psalm.}

For Bellarmine, Christ perfects the planting and governs it. Restoration is consequently more than repair of national prosperity. His conclusion distinguishes life of grace here from life of glory hereafter. The vine's final good cannot be measured simply by recovering an earlier political position or external splendor. A community can seek God's face amid diminished circumstances and still be called toward a perfection that exceeds its former success.

Bellarmine's account carries the polemical vocabulary of its own time, and his assertion that all agree on a Christological identification is his assertion, not a survey of every interpreter. The argument retained here is his specific account of continuity through apostles, converts and grafted Gentiles, with Christ as the vineyard's perfection. Augustine's corresponding account establishes genuine agreement on this point without turning the two expositions into a single voice.

Their historical interpretations also differ. Augustine understands the river as Jordan and the beast in relation to Roman devastation; Bellarmine extends the image to Euphrates and discusses Assyrian and Babylonian ravagers. These are materially different readings of the poem's geography and disaster. The shared Christian hope does not settle them. A reader can affirm their theological convergence while recognizing that neither exposition alone supplies a demonstrated date or event for every image.\footnote{Augustine, Psalm 79, §§6--8; Bellarmine, Psalm 79, explanation at vv. 10--13.}

\subsection{Restoration must answer the outcry}
The word ``restoration'' can conceal a danger: a community may want its old strength back without wanting its old injustices exposed. Isaiah prevents that evasion. The failed fruit is not smallness or vulnerability, but injustice. To restore the vineyard cannot mean giving oppressors a safer enclosure within which to continue. God's visitation must be sought as renewal of the life he planted, including a changed relation to those whose cries indict it.

Augustine makes an interior turn when explaining the poem's images of burning and digging. He distinguishes desire and fear as sources of sin. A person may be drawn into wrongdoing by an attractive reward or driven into it by a threat. He then sets good love and good fear against those distorted powers. Repentance and holy love turn the heart toward the good rather than simply preserving its former desires under religious protection.\footnote{Augustine, Psalm 79, §10. The moral exposition is Augustine's; it is distinct from the psalm's literal account of a ravaged planting.}

This moral turn deepens the corporate account. Unjust institutions are sustained by actual desires and fears: the wish to gain, the dread of losing standing, the refusal to accept correction. Matthew's tenants display both attachment to possession and hostility to the messenger who questions it. A community renewed in Christ must learn a different form of strength. It can confess wrongdoing without concluding that truth will destroy every possibility of common life, because its foundation is the rejected Son rather than its own unblemished reputation.

The Gospel's new stewardship remains stewardship. Chrysostom identifies the stone with Christ and the builders with the teachers who reject him. His exposition also describes believing Jews and Gentiles made one. These claims are more precise than a collective indictment of an ethnicity, and they leave the Christian hearer subject to the same demand for fruit. Restoration cannot transfer the tenants' possessive attitude unchanged to a new group. It must change the relation to the owner and his Son.\footnote{Chrysostom, \work{Homilies on Matthew} 68.2. The homily's broader anti-Jewish rhetoric does not license collective ethnic guilt.}

\subsection{A common life learns peace}
Philippians contributes a practical grammar of renewal. Paul has named a disagreement between Euodia and Syntyche and called for agreement in the Lord. The appointed counsel about prayer and thought then reaches a command to practice what has been learned, received, heard and seen. Peace belongs within this shared life, not outside it as a private refuge from responsibility. The community needs forms of attention and conduct that make reconciliation possible.\footnote{Philippians 4:2--9; Chrysostom, \work{Homilies on Philippians} 14.}

Chrysostom's exposition joins peace with the absence of revenge. God's gift of the Son to enemies reveals a generosity beyond human calculation. Toward the close of the homily, he calls the hearer to stop fighting and to do good rather than retaliate. His concern is not merely the avoidance of visible conflict. He names the remembrance of injuries and the motives of covetousness, envy and vainglory that feed enmity. These habits can make a community incapable of receiving the peace it professes to want.

The reading's truth and justice remain alongside that exhortation. Peace cannot require concealing harm, leaving the wronged without remedy, or praising what is false. Isaiah's outcry still sounds within the Sunday. Paul's directions toward worthy thought and action give reconciliation substance: a restored community is accountable to what is true and just. The promise of the God of peace is not a warrant to remove difficult people from view so that the surface looks calm.

John's acclamation further identifies the people's calling as lasting fruit. Augustine's account of charity makes election incompatible with superiority toward the neighbor. The disciples are chosen into love, not into a status that dispenses them from it. In his next tractate Augustine even distinguishes the creature made by God from the vice that harms it: love of enemies seeks the good of the person while refusing the evil. A restored people can therefore resist wrongdoing without making hatred its principle of unity.\footnote{Augustine, \work{Tractates on John} 87.1,4.}

\subsection{The one bread judges division}
Chrysostom's explanation of First Corinthians 10:16--17 follows the intensity of Paul's language. Participation is communion, and communion leads to the declaration that the many are one body. The preacher uses the many grains in one bread to explain how those who receive Christ are joined both to him and to one another. There is not one Christ for the hearer and a different Christ for the neighbor. Their common nourishment is the ground of their mutual obligation.\footnote{Chrysostom, \work{Homilies on First Corinthians} 24.4, NPNF1, vol. 12.}

He immediately tests that claim against conduct. Why do people who are nourished by the same gift fail to show the same love? His rebuke of actual divisions is inseparable from his strong sacramental teaching. The one body is not merely a pleasant description of a group that already behaves well. It reveals what the gift makes possible and what lovelessness contradicts. The Eucharist thus contributes a criterion of communal life more searching than shared affiliation.

Aquinas begins from the same Christological center when explaining Eucharistic grace. Christ and his Passion cause the gift; nourishment describes its sustaining effect; the species' unity from many elements signifies charity and incorporation. This prevents restoration from becoming a merely organizational project. Better arrangements may serve common life, but the unity sought in Communion has its source in Christ. The Prayer over the Offerings asks for redemption's sanctifying action, and the Prayer after Communion carries the sacramental gift toward the life of its recipients.\footnote{Aquinas, \work{Summa theologiae} III, q. 79, a. 1, body and reply 2.}

The Corinthians antiphon is an appointed alternative, however, not the compulsory ending of this Sunday. If the Lamentations antiphon is used, the immediate accent is seeking and waiting. The wounded people may still be unable to see the repair for which they long. This option protects hope from becoming a demand that pain promptly resolve into a harmonious picture. Both alternatives belong to Eucharistic participation, but only one makes the one-body language explicit at that moment.

\subsection{Waiting without surrendering hope}
Lamentations speaks hope amid a sustained account of suffering. The speaker recalls bitterness, yet turns toward the Lord's mercies. The chosen verse joins goodness with seeking and patient expectation. It does not pretend that the city's wound has vanished. The full chapter's protest against injustice and its later appeal for vindication forbid interpreting patience as consent to whatever a stronger person does. Hope seeks the good Lord, not the preservation of oppression.\footnote{Lamentations 3:17--36,40--58.}

For a wounded community this is a demanding form of truthfulness. It can ask for life before it knows how restored life will look. It can refuse to treat the present ruin as the whole meaning of its history. It can also repent without claiming that every injured member is equally responsible for the wrong. The psalm's corporate petition makes room for a common plea while Isaiah and the Gospel retain their particular accusations. Communion nourishes this people in its need, not a fictional people whose conflicts are already resolved.

The final prayer's transformation therefore includes a task as well as consolation. Leo's Eucharistic account describes receivers carrying the crucified and risen Christ in body and spirit. Aquinas's reply about the body gives this an immediate and future horizon: bodily members serve justice in the present, and the body shares incorruption in the life to come. The restored vineyard is called into embodied charity now while still awaiting its consummation.\footnote{Leo, Sermon 63, VII; Aquinas, III, q. 79, a. 1, reply 3.}

\subsection{The reward is more than recovery}
Augustine's psalm exposition closes by refusing to make created benefits the final reason for worship. God himself is the reward. Bellarmine distinguishes grace from glory. Their convergence keeps restoration open toward a good that no return to former security can exhaust. A people may be renewed in faith and charity while still carrying scars, losses and unresolved grief. Its hope rests in God rather than in a guaranteed restoration of every earthly arrangement.

This also places a limit upon triumph. The Church cannot identify its present condition without remainder with the completed vineyard. Matthew's fruit criterion continues to judge its stewardship; Isaiah's justice continues to expose its failures; Paul's practiced peace continues to require action. The promise is neither that the wounded will remain abandoned nor that the renewed are beyond correction. Christ gathers, nourishes and perfects a people who still pray for his face to shine upon them.

\begin{fourSenses}
\item[Literal.] Israel's planting is accused and its broken condition lamented; Jesus confronts the stewards who resist God's messengers; Paul calls a real community toward prayer and practiced peace. The voices of guilt, injury and hope remain distinct within common life.
\item[Allegorical.] Augustine and Bellarmine read the vine's perfection in Christ, the rejected Son and cornerstone. Jewish apostles and converts, together with grafted Gentiles, express continuity in the renewed people; the one bread joins the many to Christ and one another.
\item[Moral.] Renewal requires justice, truthful repentance, reconciliation and service of those who share the gift. Unity cannot conceal wrongdoing or silence its victims. Patient hope can name pain while refusing hatred and despair.
\item[Anagogical.] Grace tends toward glory, God's visitation toward unhindered fellowship, and sacramental nourishment toward bodily resurrection. The final good is God himself; no earthly community's prosperity exhausts that promise.
\end{fourSenses}
